\chapter{牛顿运动定律}

牛顿运动定律把运动学中的“如何描述运动”推进到“为什么会这样运动”。
在前一章中，我们已经能够用位置 $x(t)$、速度 $v(t)=\dv{x}{t}$、加速度 $a(t)=\dv[2]{x}{t}$ 描述一个质点的运动；
而在本章中，我们将建立力与加速度之间的定量联系，并把高中力学中大量熟悉的题型统一写成微分方程初值问题：
给定受力 $\vb{F}$ 与初始条件，求解质点的轨迹 $\vb{r}(t)$、速度 $\vb{v}(t)$ 与位移规律。

从微积分的角度看，牛顿力学最核心的一句话可以写成
\[
  m\dv[2]{\vb{r}}{t}=\vb{F}(\vb{r},\dv{\vb{r}}{t},t).
\]
它告诉我们：力学不是若干零散公式的堆砌，而是一套围绕二阶常微分方程展开的理论。

\section{牛顿第一定律与惯性}
\subsection{第一定律的内容}

\begin{law}[牛顿第一定律]
若物体不受外力，或所受合外力为零，则它将保持静止状态或做匀速直线运动。
\end{law}

这一表述表面上像是在讨论“没有力时会怎样”，实质上却是在回答一个更根本的问题：
\emph{什么样的参考系才适合用最简单的力学规律来描述世界？}

\begin{definition}[惯性参考系]
在某一参考系中，若孤立质点总保持静止或匀速直线运动，则称该参考系为\textbf{惯性参考系}。
\end{definition}

地面参考系在许多高中问题中可以近似看作惯性系；
而相对于地面做匀速直线运动的列车车厢，也可以近似视为惯性系。
这说明惯性系并不唯一，它们彼此之间做匀速直线运动。

\begin{theorem}[惯性系之间的相对运动]
若参考系 $S$ 是惯性系，而参考系 $S'$ 相对 $S$ 做匀速直线运动，则 $S'$ 也是惯性系。
\end{theorem}

\begin{proof}
设质点在 $S$ 中不受外力，因此有
\[
  \dv[2]{\vb{r}}{t}=\vb{0}.
\]
若 $S'$ 相对 $S$ 的速度为常矢量 $\vb{u}$，则坐标变换为
\[
  \vb{r}'=\vb{r}-\vb{u}t-\vb{r}_0.
\]
对时间求两次导数得
\[
  \dv[2]{\vb{r}'}{t}=\dv[2]{\vb{r}}{t}=\vb{0}.
\]
因此在 $S'$ 中孤立质点仍做静止或匀速直线运动，所以 $S'$ 也是惯性系。
\end{proof}

\begin{remark}
第一定律不是第二定律在 $\vb{F}=0$ 时的“特例”，因为它先给出了惯性系的判据，
为第二定律提供了适用背景。没有惯性系的概念，$\vb{F}=m\vb{a}$ 本身没有明确物理意义。
\end{remark}

\subsection{惯性的物理意义}

\begin{definition}[惯性]
物体保持原有运动状态不变的性质称为\textbf{惯性}。质量越大，改变其速度所需的作用越困难，因而质量是惯性大小的量度。
\end{definition}

我们把“运动状态”理解为速度 $\vb{v}$。若想让 $\vb{v}$ 改变，就必须让
\[
  \dv{\vb{v}}{t}\neq \vb{0},
\]
也就是必须产生加速度。于是，力学问题自然过渡到“力与加速度的关系”。

\paragraph{惯性的直观例子：匀速列车上抛小球}
以地面为参考系，小球出手瞬间除了竖直速度外，还具有与列车相同的水平速度 $v_0$。
若忽略空气阻力，其水平方向合力为零，因此
\[
  m\dv[2]{x}{t}=0 \quad \Longrightarrow \quad \dv{x}{t}=v_0,
\]
于是小球在空中始终与车厢保持相同的水平位移；
竖直方向仅受重力作用，因此做上抛运动。故它仍落回抛出者附近。

这正是惯性最直观的表现：物体不会因为“离开手”就失去原有的水平速度。

\subsection{由第一定律到受力分析}

在实际问题中，我们往往不直接判断“物体为什么动”，而是先判断它\emph{受到哪些力}。
若合力为零，则加速度为零，速度保持不变；若合力不为零，则速度必然改变。
因此，牛顿第一定律告诉我们：

\begin{center}
  \fbox{\parbox{0.8\textwidth}{
  判断一个力学问题的第一步，不是代公式，而是辨认参考系是否近似惯性，并判断合力是否为零。
  }}
\end{center}

\section{牛顿第二定律}
\subsection{定律及其微分方程形式}

\begin{law}[牛顿第二定律]
在惯性参考系中，质点所受合外力等于质量与加速度的乘积，即
\[
  \vb{F}_{\text{合}}=m\vb{a}.
\]
当质量恒定时，也可写为
\[
  \vb{F}_{\text{合}}=m\dv{\vb{v}}{t}=m\dv[2]{\vb{r}}{t}.
\]
\end{law}

若只研究一维直线运动，则第二定律写成
\[
  m\dv[2]{x}{t}=F(x,\dot{x},t),
\]
其中 $\dot{x}=\dv{x}{t}$。这就是一个二阶常微分方程。
要确定唯一运动，还必须配上初始条件
\[
  x(0)=x_0,\qquad \dot{x}(0)=v_0.
\]
因此一个完整的动力学问题，本质上是一个\textbf{二阶初值问题}。

\begin{theorem}[初值问题的决定性]
若一维质点满足
\[
  m\dv[2]{x}{t}=F(x,\dot{x},t),
\]
且函数 $F$ 对变量足够连续，则给定初始位置 $x(0)=x_0$ 和初速度 $\dot{x}(0)=v_0$ 后，运动在局部范围内唯一确定。
\end{theorem}

这个定理的数学证明属于常微分方程理论，但物理意义非常重要：
\emph{经典力学是可预言的。}只要受力规律和初态已知，未来运动就由方程唯一决定。

\subsection{标量形式与分量形式}

在平面或空间中，向量方程
\[
  m\dv[2]{\vb{r}}{t}=\vb{F}
\]
可以沿坐标轴分解为
\[
  m\dv[2]{x}{t}=F_x,\qquad m\dv[2]{y}{t}=F_y,\qquad m\dv[2]{z}{t}=F_z.
\]
这说明二维、三维问题往往可以分解为若干一维微分方程分别求解。
抛体运动之所以容易处理，根本原因就在于水平方向与竖直方向的方程可以分离。

\subsection{受力图与动力学建模}

写出微分方程之前，最重要的中间步骤是受力分析。
典型流程为：
\begin{enumerate}
  \item 选定研究对象；
  \item 选定参考系与坐标轴；
  \item 画出受力图；
  \item 沿坐标方向列出牛顿第二定律；
  \item 加入约束关系或初始条件；
  \item 解微分方程并检查物理意义。
\end{enumerate}

\begin{figure}[H]
  \centering
  \begin{tikzpicture}[scale=1.0, >=Latex]
    \draw[thick] (0,0) -- (5,0);
    \fill[gray!20] (2,0) rectangle (3.4,1);
    \node at (2.7,0.5) {$m$};
    \draw[->, thick, blue] (2.7,1) -- (2.7,2) node[right] {$\vb{N}$};
    \draw[->, thick, red] (2.7,0.5) -- (2.7,-1.3) node[right] {$m\vb{g}$};
    \draw[->, thick, orange!80!black] (3.4,0.5) -- (4.8,0.5) node[above] {$\vb{F}$};
    \draw[->, thick, teal!70!black] (2,0.5) -- (0.7,0.5) node[above] {$\vb{f}$};
  \end{tikzpicture}
  \caption{水平面上物块的受力图：支持力、重力、外推力与摩擦力}
\end{figure}

图中的水平方向方程为
\[
  m\dv[2]{x}{t}=F-f,
\]
竖直方向若无加速度，则有 $N-mg=0$。若把摩擦力大小随外加力变化的过程画成图像，可以更直观地看见：静摩擦力先随外力同步增大，达到最大静摩擦后才转入较小的动摩擦力。

\begin{figure}[htbp]
\centering
\begin{tikzpicture}
\begin{axis}[
  width=0.75\textwidth, height=0.5\textwidth,
  xlabel={外加力 $F_{\text{施}}$}, ylabel={摩擦力 $F_{\text{摩}}$},
  axis lines=middle,
  grid=major, grid style={gray!30},
  thick,
  xmin=0, xmax=4.2, ymin=0, ymax=1.4,
]
\addplot[blue] coordinates {(0,0) (1,1)};
\addplot[blue] coordinates {(1,0.7) (4,0.7)};
\addplot[only marks, mark=*, blue] coordinates {(1,0.7)};
\addplot[only marks, mark=o, blue] coordinates {(1,1)};
\addplot[dashed] coordinates {(1,0) (1,1.05)};
\node at (axis cs:0.55,1.12) {静摩擦};
\node at (axis cs:2.6,0.88) {动摩擦};
\end{axis}
\end{tikzpicture}
\caption{摩擦力随外加力的变化：静摩擦先“跟着走”，滑动后转为近似恒定的动摩擦}
\end{figure}

许多经典题的“难点”，其实只是把这些步骤压缩得过于简短。

\begin{example}[恒力作用下的直线运动]
质量为 $m$ 的小车在水平光滑轨道上受恒力 $F_0$ 作用，初始时刻位置为 $x_0$，速度为 $v_0$。求其运动方程。
\end{example}

\begin{proof}
由第二定律，
\[
  m\dv[2]{x}{t}=F_0.
\]
两边积分一次，得
\[
  \dv{x}{t}=\frac{F_0}{m}t+C_1.
\]
由初速度条件 $\dot{x}(0)=v_0$，知 $C_1=v_0$，故
\[
  v(t)=v_0+\frac{F_0}{m}t.
\]
再积分一次，
\[
  x(t)=x_0+v_0 t+\frac{F_0}{2m}t^2.
\]
这就是匀变速直线运动公式的微积分来源。
\end{proof}

\begin{figure}[htbp]
\centering
\begin{tikzpicture}
\begin{axis}[
  width=0.75\textwidth, height=0.5\textwidth,
  xlabel={质量 $m$}, ylabel={加速度 $a$},
  axis lines=middle,
  grid=major, grid style={gray!30},
  thick,
  xmin=0, xmax=5.5, ymin=0, ymax=4.5,
]
\addplot[blue, domain=0.5:5, samples=100] {2/x};
\end{axis}
\end{tikzpicture}
\caption{在合力恒定时，$a=F/m$ 表明质量越大，加速度越小}
\end{figure}

\subsection{质量与加速度的方向}

牛顿第二定律是一个向量定律，加速度方向总与合力方向一致，而不一定与速度方向一致。
例如竖直上抛运动中，物体上升时速度向上，但合力向下，因此加速度向下；
正是这种“速度与加速度可反向”的情形，导致速度减小但运动仍继续向上。

\begin{remark}
在解题时，经常把“合力方向”“速度方向”“位移方向”混为一谈。微积分语言可以帮助我们澄清：
速度是 $\dv{x}{t}$ 的符号，加速度是 $\dv[2]{x}{t}$ 的符号，而位移只是 $x$ 的变化。
三者相关，但并不等同。
\end{remark}

\section{牛顿第三定律}
\subsection{作用与反作用}

\begin{law}[牛顿第三定律]
两个物体之间的相互作用力总是大小相等、方向相反，并且作用在同一直线上。
\end{law}

若物体 $A$ 对物体 $B$ 施力 $\vb{F}_{AB}$，则物体 $B$ 同时对物体 $A$ 施力 $\vb{F}_{BA}$，满足
\[
  \vb{F}_{AB}=-\vb{F}_{BA}.
\]

第三定律常被误解为“相互抵消”，但需要特别注意：
作用力与反作用力分别作用在不同物体上，因此不能在同一个受力图中相加为零。

\begin{definition}[相互作用力对]
由两个物体相互作用产生、大小相等方向相反、分别作用于两个物体上的一对力，称为\textbf{作用力与反作用力}。
\end{definition}

\subsection{第三定律与系统观点}

对于由两个物体组成的系统，若只考虑它们之间的内力，则
\[
  \vb{F}_{12}+\vb{F}_{21}=\vb{0}.
\]
这为后续学习动量守恒奠定基础：系统内部相互作用不会改变总动量，只能改变内部各部分的运动状态。

\paragraph{生活中的第三定律：人推墙时为什么手会疼}
人的手对墙施加推力 $\vb{F}_{\text{手对墙}}$，同时墙对手施加反作用力
\[
  \vb{F}_{\text{墙对手}}=-\vb{F}_{\text{手对墙}}.
\]
墙之所以几乎不动，是因为它与地基构成巨大系统，质量远大，产生的加速度极小；
而手部组织柔软，更容易形变，于是人会感到疼痛。第三定律说明了力的配对关系，但运动效果仍须结合第二定律和质量大小一起分析。

\begin{remark}
常见错误说法有两种：
\begin{enumerate}
  \item “马拉车时，马对车的拉力和车对马的拉力相等，所以车不会动”；
  \item “书放在桌面上，书受重力与桌面对书的支持力是一对作用力与反作用力”。
\end{enumerate}
第一种忽略了这两个力作用在不同物体上；
第二种中，重力与支持力作用在同一物体上，它们是一对平衡力，而不是作用力与反作用力。
\end{remark}

\section{常见力学问题的微分方程解法}

本节把若干熟悉的高中模型统一到微分方程框架中。请注意：
\emph{真正的核心不是记住答案，而是学会从受力到方程、再从方程到运动的全过程。}

\subsection{恒力模型}

若质点在一维上受恒力 $F_0$，则满足
\[
  m\dv[2]{x}{t}=F_0.
\]
其通解为
\[
  x(t)=C_1+C_2 t+\frac{F_0}{2m}t^2.
\]
加入初值条件后，就得到熟悉的匀变速运动规律。

特别地，在竖直方向取向上为正，则自由落体满足
\[
  m\dv[2]{y}{t}=-mg,
\]
即
\[
  \dv[2]{y}{t}=-g.
\]
这说明自由落体不是“因为速度越来越大而加速”，而是因为恒定重力使加速度近似恒定。

\begin{example}[竖直上抛的微分方程解]
物体以初速度 $v_0$ 从地面竖直上抛，不计空气阻力。求速度、位移与最高点时间。
\end{example}

\begin{proof}
取竖直向上为正，运动方程为
\[
  m\dv{v}{t}=-mg,
\]
即
\[
  \dv{v}{t}=-g.
\]
积分得
\[
  v(t)=v_0-gt.
\]
再积分并用 $y(0)=0$，得
\[
  y(t)=v_0 t-\frac{1}{2}gt^2.
\]
最高点满足 $v(t_\text{max})=0$，故
\[
  t_\text{max}=\frac{v_0}{g},
\]
代回得最大高度
\[
  y_{\text{max}}=\frac{v_0^2}{2g}.
\]
熟悉的结论由微分方程自然推出。
\end{proof}

\subsection{弹簧力：$F=-kx$}

理想弹簧满足胡克定律
\[
  F=-kx,
\]
其中 $x$ 是相对平衡位置的位移，$k>0$ 是劲度系数。于是运动方程为
\[
  m\dv[2]{x}{t}=-kx,
\]
即
\begin{equation}
  \dv[2]{x}{t}+\omega^2 x=0,
  \qquad \omega\coloneqq\sqrt{\frac{k}{m}}.
  \label{eq:sho}
\end{equation}
这就是最基本的简谐振动方程。

其通解可写为
\[
  x(t)=A\cos(\omega t)+B\sin(\omega t),
\]
也可写成
\[
  x(t)=C\cos(\omega t+\varphi).
\]

\begin{figure}[H]
  \centering
  \begin{tikzpicture}[scale=1.0, >=Latex]
    \draw[fill=gray!20] (-0.5,-0.8) rectangle (0,0.8);
    \draw[thick, decorate, decoration={coil, aspect=0.45, segment length=4pt, amplitude=4pt}] (0,0) -- (4,0);
    \fill[blue!15] (4,-0.7) rectangle (5.4,0.7);
    \node at (4.7,0) {$m$};
    \draw[dashed] (4.7,-1.2) -- (4.7,1.2);
    \draw[->, thick] (4.7,-1.0) -- (6.1,-1.0) node[below] {$x$};
    \draw[->, thick, red] (4.7,0) -- (3.2,0) node[above] {$-k x$};
  \end{tikzpicture}
  \caption{水平弹簧振子：恢复力总指向平衡位置}
\end{figure}

\begin{example}[弹簧振子的初值问题]
水平光滑轨道上，质量为 $m$ 的物块与弹簧相连。以平衡位置为原点，初始时刻有 $x(0)=x_0$，$v(0)=0$。求运动规律。
\end{example}

\begin{proof}
由 \cref{eq:sho}，通解写为
\[
  x(t)=A\cos(\omega t)+B\sin(\omega t),\qquad \omega=\sqrt{\frac{k}{m}}.
\]
由初值 $x(0)=x_0$，得 $A=x_0$；
再求导得
\[
  v(t)=\dv{x}{t}=-A\omega\sin(\omega t)+B\omega\cos(\omega t).
\]
由 $v(0)=0$ 得 $B=0$。因此
\[
  x(t)=x_0\cos\left(\sqrt{\frac{k}{m}}\,t\right),
\]
\[
  v(t)=-x_0\sqrt{\frac{k}{m}}\sin\left(\sqrt{\frac{k}{m}}\,t\right).
\]
可见位移与速度都是周期函数，其周期为
\[
  T=2\pi\sqrt{\frac{m}{k}}.
\]
\end{proof}

\subsection{阻力模型：$F=-bv$}

在速度不太大时，阻力常近似与速度成正比，写为
\[
  F_{\text{阻}}=-bv,
\]
其中 $b>0$。若取竖直向下为正，物体在空气中下落时的方程为
\[
  m\dv{v}{t}=mg-bv.
\]
这是一个一阶线性微分方程。

整理为
\[
  \dv{v}{t}+\frac{b}{m}v=g.
\]
其解为
\[
  v(t)=\frac{mg}{b}+C\exp\left(-\frac{b}{m}t\right).
\]
若 $v(0)=0$，则
\begin{equation}
  v(t)=\frac{mg}{b}\left(1-\exp\left(-\frac{b}{m}t\right)\right).
  \label{eq:dragv}
\end{equation}
当 $t\to\infty$ 时，速度趋于
\[
  v_T=\frac{mg}{b},
\]
称为\textbf{终端速度}。

\begin{example}[线性阻力下的下落]
质量为 $m$ 的小球自静止开始竖直下落，受重力和线性阻力 $-bv$。求速度随时间的变化，并解释其物理意义。
\end{example}

\begin{proof}
方程与上文一致，其解即为 \cref{eq:dragv}。
由于指数项逐渐衰减，速度不是无限增大，而是逐渐逼近终端速度 $v_T=mg/b$。
当 $v$ 很小时，阻力小，故加速度接近 $g$；
当 $v$ 接近 $v_T$ 时，$mg-bv\to 0$，加速度趋近于零，速度便稳定下来。

这体现了微分方程语言的优势：它不仅给出“最后怎样”，还描述了“怎样逐渐逼近”。
\end{proof}

\subsection{综合问题：重力、弹簧力与阻力同时存在}

若把一维向下位移记作 $x$，一个物体在竖直方向同时受到重力、弹簧力与线性阻力，则可写出
\[
  m\dv[2]{x}{t}=mg-kx-b\dv{x}{t}.
\]
整理得
\begin{equation}
  m\dv[2]{x}{t}+b\dv{x}{t}+kx=mg.
  \label{eq:damped-shifted}
\end{equation}
这是高中常见“弹簧+阻尼+重力”综合模型的标准方程。

令
\[
  x_e=\frac{mg}{k},
\]
则 $x_e$ 表示静止平衡位置。引入新变量
\[
  y=x-x_e,
\]
方程化为
\[
  m\dv[2]{y}{t}+b\dv{y}{t}+ky=0.
\]
这说明重力的作用主要是\emph{平移平衡位置}，而围绕平衡位置的小幅运动仍由齐次方程控制。

\begin{example}[带阻尼的竖直弹簧]
某物块悬挂在轻弹簧下端，释放后在空气中做小幅振动。设向下为正，以弹簧原长位置为 $x=0$，阻力满足 $F=-b\dot{x}$。说明如何求解其运动。
\end{example}

\begin{proof}
受力有三项：重力 $mg$ 向下，弹簧力 $-kx$ 向上，阻力 $-b\dot{x}$ 总与速度反向。故有
\[
  m\dv[2]{x}{t}+b\dv{x}{t}+kx=mg.
\]
先求静止平衡位置 $x_e=mg/k$，再令 $y=x-x_e$，得到齐次方程
\[
  m\dv[2]{y}{t}+b\dv{y}{t}+ky=0.
\]
设特征方程为
\[
  mr^2+br+k=0.
\]
若 $b^2<4mk$，则根为
\[
  r=-\frac{b}{2m}\pm i\sqrt{\frac{k}{m}-\frac{b^2}{4m^2}},
\]
于是解为
\[
  y(t)=e^{-\frac{b}{2m}t}\left(C_1\cos \omega_d t+C_2\sin \omega_d t\right),
\]
其中
\[
  \omega_d=\sqrt{\frac{k}{m}-\frac{b^2}{4m^2}}.
\]
最后再写回 $x(t)=x_e+y(t)$ 即可。它表示物体围绕新的平衡位置做振幅逐渐减小的振动。
\end{proof}

\subsection{从模型到方法}

本节几种模型看似不同，实则都服从同一流程：
\[
  \text{受力分析}\ \longrightarrow\ \text{微分方程}\ \longrightarrow\ \text{初值条件}\ \longrightarrow\ \text{解析解或定性分析}.
\]

当方程能直接积分时，我们得到显式公式；
当方程较复杂时，即使写不出闭式解，也可通过平衡位置、极限行为、单调性等进行定性讨论。
这正是高中力学与高等数学自然衔接的地方。

\section{非惯性系与惯性力}
\subsection{加速参考系中的运动方程}

前面所有定律都默认参考系是惯性系。若参考系本身具有加速度，事情就会发生变化。
设惯性系 $S$ 中坐标为 $x$，另有参考系 $S'$ 相对 $S$ 沿 $x$ 方向以加速度 $a_0$ 运动。
若两系原点在 $t=0$ 重合，且初始相对速度为零，则有
\[
  x=x'+\frac{1}{2}a_0 t^2.
\]
两次求导得
\[
  \dv[2]{x}{t}=\dv[2]{x'}{t}+a_0.
\]
在惯性系 $S$ 中，物体满足
\[
  m\dv[2]{x}{t}=F.
\]
代入上式得
\[
  m\dv[2]{x'}{t}=F-ma_0.
\]

\begin{definition}[惯性力]
在加速参考系中，为了保持牛顿第二定律形式不变而引入的附加力
\[
  \vb{F}_{\text{惯}}=-m\vb{a}_0
\]
称为\textbf{惯性力}。
\end{definition}

因此在非惯性系中可以写成
\[
  m\vb{a}'=\vb{F}_{\text{真实}}+\vb{F}_{\text{惯}}.
\]
这里的惯性力不是由实际相互作用产生，而是参考系加速带来的表观效应。

\begin{example}[加速电梯中的“超重”与“失重”]
质量为 $m$ 的人站在电梯中。若电梯以加速度 $a_0$ 向上运动，求地板对人的支持力，并解释超重现象。
\end{example}

\begin{proof}
在地面近似惯性系中，取竖直向上为正。人受支持力 $N$ 向上、重力 $mg$ 向下，故
\[
  N-mg=ma_0.
\]
于是
\[
  N=m(g+a_0).
\]
当 $a_0>0$ 时，支持力大于重力，人感觉“更重”；
当电梯向下加速，即 $a_0<0$ 时，支持力减小，出现“失重感”。

若改在电梯参考系中看，由于参考系向上加速，应加入向下的惯性力 $ma_0$，则平衡写作
\[
  N-mg-ma_0=0,
\]
结论完全一致。
\end{proof}

\subsection{初步理解等效原理}

在局部范围内，一个均匀加速参考系中的惯性力，与一个均匀重力场在数学形式上十分相似。
例如在加速度为 $a_0$ 的电梯中，物体所受“有效重力”可写作
\[
  g_{\text{eff}}=g+a_0
\]
或 $g-a_0$，取决于方向约定。于是许多现象——摆的周期、弹簧秤读数、液面形状——都可理解为对有效重力的响应。

\begin{remark}
爱因斯坦正是从这种“引力场与加速参考系局部等效”的思想出发，逐步发展出广义相对论。
在本书当前层次中，我们只需把它理解为：
\emph{某些非惯性效应可以等价地看成一种附加的场。}
\end{remark}

\subsection{为什么转弯时会“被甩出去”}

乘车转弯时，人会感觉自己被“甩向外侧”。在地面惯性系中，这并不是有一个真实的外甩力，
而是人体由于惯性倾向于保持原来的直线运动；
车门或座椅提供向内的力，迫使人体改变速度方向。

而在随车转动的非惯性系中，为了保持形式上的牛顿第二定律，必须引入离心惯性力等项。
因此，“被甩出去”的感觉本质上是\emph{非惯性系描述中的惯性力表现}。

下面选取三类极具代表性的经典力学题，用微分方程语言重写。这样做的目的不是把题目“做复杂”，
而是帮助我们看见：高中题型背后有统一的动力学结构。

\section*{本章小结}
\addcontentsline{toc}{section}{本章小结}

\begin{enumerate}
  \item 牛顿第一定律给出惯性系的判据，说明孤立物体保持静止或匀速直线运动。
  \item 牛顿第二定律在一维中可写成
  \[
    m\dv[2]{x}{t}=F(x,\dot{x},t),
  \]
  它是经典力学最核心的微分方程。
  \item 牛顿第三定律揭示相互作用力成对出现，但作用在不同物体上。
  \item 恒力、弹簧力、阻力以及综合模型，都可以通过建立初值问题系统求解。
  \item 在非惯性系中需引入惯性力；加速参考系与重力场的局部等效思想，是理解相对论的重要起点之一。
\end{enumerate}

\addcontentsline{toc}{section}{习题}

